%======================================================================
\chapter{The folding theorem}\label{ch:folding}
%======================================================================

This chapter proves the theorem about crooked chains that Part~II of the uniqueness proof needs: a
chain crooked in $D$, a stretch of which reaches both ends of $D$, can be consolidated into a chain
whose first and last blocks are prescribed and whose blocks are small. It is Bing's Theorem~4 (1948),
in the form that the formal proof establishes.

Throughout, $D$ is a chain of length $k$, $E$ is a chain of length $n$ which is crooked in $D$, and
$\lo(t)$, $\hi(t)$ are the least and the greatest index $j$ with $E_t\subseteq D_j$
(Definition~\ref{def:minidx}). For integers $w$ we write $D[w,w+2]$ for the union of those of
$D_w,D_{w+1},D_{w+2}$ that exist. A block map (Definition~\ref{def:block-map}) has \emph{small
blocks} if every block $E^\beta_u$ lies in some $D[w,w+2]$.

\section{The statement}

\begin{theorem}[Folding theorem; Bing 1948, Theorem~4]\label{thm:folding}
Let $k\ge3$, let $E$ be crooked in $D$, and let $r,r',s',s$ be indices of $E$ with $r'=r+1$, $s=s'+1$ and
$r'\le s'$. Suppose that some $E_h$ with $r\le h\le s$ meets $D_0$, and some $E_{h'}$ with $r\le h'\le s$
meets $D_{k-1}$. Then there is a block map $\beta\colon[n]\to[q]$ with small blocks such that
\[
\beta(r)=\beta(r')=0\qquad\text{and}\qquad\beta(s')=\beta(s)=q-1 .
\]
\end{theorem}

\formal{\lean{HPC.Fold.theorem4_pair} (in \lean{HPC/Fold/Theorem4Pair.lean}). The version with single links
$r$ and $s$ in place of the pairs is \lean{HPC.Fold.theorem4}.}

We need the pairs because of the following application. A point lies in at most two links of a chain, and
these are adjacent; for the block chain to go from $R$ to $S$ (Lemma~\ref{lem:block-from}) \emph{every} link
containing $R$ must be in the first block.

\begin{corollary}\label{cor:block-chain}
Let $k\ge5$ and let $E$ be crooked in $D$. Let $a<b$ be indices of $E$ such that the subchain from $a$ to $b$
goes from $R$ to $S$, and suppose some $E_{t_1}$ and some $E_{t_2}$ with $a\le t_1,t_2\le b$ meet $D_0$ and $D_{k-1}$
respectively. Then there is a block map $\beta$ for $E$ with small blocks such that the block chain $E^\beta$ goes
from $R$ to $S$.
\end{corollary}

\begin{proof}
First, $b\ge a+2$: the links $E_{t_1}$ and $E_{t_2}$ lie in links of $D$ with indices $\le1$ and $\ge k-2\ge3$
respectively, so they do not meet, $\abs{t_1-t_2}\ge2$. Next, the links of $E$ containing $R$: $R$ lies in $E_a$, and in
no $E_t$ with $a<t\le b$; since links containing $R$ meet, they are among $E_{a-1},E_a$. Similarly the links containing
$S$ are among $E_b,E_{b+1}$. Apply Theorem~\ref{thm:folding} with $r=a-1$, $r'=a$ if $a\ge1$ (and $r=a$, $r'=a+1$ if
$a=0$), and $s'=b$, $s=b+1$ if $b+1<n$ (and $s'=b-1$, $s=b$ if $b=n-1$). The hypotheses hold since $r\le a\le t_1,t_2\le
b\le s$ and $r'\le a+1\le b-1\le s'$. Every link containing $R$ is $E_r$ or $E_{r'}$, hence in block $0$, and every link
containing $S$ is in the last block; by Lemma~\ref{lem:block-from}, $E^\beta$ goes from $R$ to $S$.
\end{proof}

\formal{\lean{HPC.Chain.exists_blockChain} (in \lean{HPC/PseudoArc/T13Step.lean}).}

\paragraph{Plan of the proof.}
The whole proof is about the integer sequences $\lo$ and $\hi$. Section~\ref{sec:fold-readings} records their
properties: they move by at most one, and $\lo$ is \emph{crooked} in a sense that only involves integers. Sections
\ref{sec:fold-folds}--\ref{sec:fold-zigzag} prove the central result, the \emph{zigzag theorem}: from any base link, a
crooked integer sequence can be ``folded'' onto a walk that goes back and forth, with every fibre of the folding map
inside three consecutive levels. This already gives the version of the theorem with only one prescribed block (Bing's
Theorem~3). To prescribe both the first and the last block, Section~\ref{sec:fold-fake} folds an initial segment of $E$
from $r$ and a final segment from $s$, and Section~\ref{sec:fold-splice} splices the two folds together with a grouping
of the middle part by levels. Section~\ref{sec:fold-proof} finishes the proof.

\section{Readings and crooked sequences}\label{sec:fold-readings}

\begin{lemma}\label{lem:readings}
For all $t$ (and $t+1$ when it is an index):
\begin{enumerate}[label=(\alph*)]
\item $\lo(t)\le\hi(t)\le\lo(t)+1$, and $E_t\subseteq D_j$ exactly for $j\in\set{\lo(t),\hi(t)}$;
\item $\abs{\lo(t+1)-\lo(t)}\le1$, $\abs{\hi(t+1)-\hi(t)}\le1$, $\abs{\hi(t+1)-\lo(t)}\le1$ and $\abs{\hi(t)-\lo(t+1)}\le 1$;
\item if $E_t$ meets $D_j$, then $\abs{\lo(t)-j}\le1$ and $\abs{\hi(t)-j}\le1$;
\item if $\lo(t)\le j\le\hi(t)$, then $E_t\subseteq D_j$.
\end{enumerate}
\end{lemma}

\begin{proof}
(a), (d) and the first two parts of (b) are Lemma~\ref{lem:minidx}. For the others: $E_t$ and $E_{t+1}$ meet at a point
$z$, which lies in $D_{\hi(t+1)}$ and in $D_{\lo(t)}$, so these links meet; likewise $D_{\hi(t)}$ and $D_{\lo(t+1)}$. (c)
$E_t$ lies in $D_{\lo(t)}$ and in $D_{\hi(t)}$ and meets $D_j$.
\end{proof}

\formal{\lean{HPC.Fold.IsCrookedIvF} and \lean{HPC.Fold.isCrookedIvF_minIdx_maxIdx} package these facts with
crookedness; see also \lean{HPC.Fold.B79.seq_facts}.}

\begin{definition}[Crooked sequence]\label{def:crooked-seq}
Let $I=[\ell,h]$ be an interval of integers. A function $f\colon I\to\Z$ with unit steps is \emph{crooked} if for all
$x<y$ in $I$:
\begin{enumerate}[label=(\roman*)]
\item if $f(y)\ge f(x)+3$, there are $x<u<u'<y$ with $f(y)-2\le f(u)\le f(y)-1$ and $f(x)\le f(u')\le f(x)+1$;
\item if $f(y)\le f(x)-3$, there are $x<u<u'<y$ with $f(y)\le f(u)\le f(y)+1$ and $f(x)-2\le f(u')\le f(x)-1$.
\end{enumerate}
\end{definition}

In words: between two places whose values differ by three or more, the sequence first comes close to the later value
and then close to the earlier one. (The asymmetry between (i) and (ii) reflects the fact that $\lo$ is the
\emph{least} index.)

\formal{\lean{HPC.Fold.Z77.CrkOn}.}

\begin{lemma}\label{lem:crooked-seq}
\begin{enumerate}[label=(\alph*)]
\item $\lo$ is crooked on $[0,n-1]$.
\item $\bar\lo:=k-1-\hi$ is crooked on $[0,n-1]$; it is the least-index reading of $E$ in the reversed chain $\overline D$.
\item If $f$ is crooked on $[\ell,h]$, then $x\mapsto f(-x)$ is crooked on $[-h,-\ell]$, and $x\mapsto f(x+c)$ is crooked on
$[\ell-c,h-c]$.
\end{enumerate}
\end{lemma}

\begin{proof}
(a) Let $x<y$ with $\lo(y)\ge\lo(x)+3$. The link $E_x$ meets $D_{\lo(x)}$ and $E_y$ meets $D_{\lo(y)}$, so crookedness of $E$ in $D$
gives $u,u'$ with $\mathrm{Btw}(x,u,u',y)$, that is $x<u<u'<y$, and $E_u\subseteq D_{\lo(y)-1}$, $E_{u'}\subseteq D_{\lo(x)+1}$.
By Lemma~\ref{lem:readings}(a), $\lo(y)-1\in\set{\lo(u),\lo(u)+1}$ and $\lo(x)+1\in\set{\lo(u'),\lo(u')+1}$, which is~(i).
If $\lo(y)\le\lo(x)-3$, crookedness applied to the pair $E_y$ (meeting $D_{\lo(y)}$) and $E_x$ (meeting $D_{\lo(x)}$) gives
$\mathrm{Btw}(y,v,v',x)$, i.e.\ $x<v'<v<y$, with $E_v\subseteq D_{\lo(x)-1}$ and $E_{v'}\subseteq D_{\lo(y)+1}$. Take $u=v'$ and
$u'=v$: then $\lo(u)\in\set{\lo(y),\lo(y)+1}$ and $\lo(u')\in\set{\lo(x)-2,\lo(x)-1}$, which is~(ii).

(b) The link $E_t$ lies in $\overline D_j=D_{k-1-j}$ exactly for $k-1-j\in\set{\lo(t),\hi(t)}$, so its least index in $\overline D$
is $k-1-\hi(t)$. $E$ is crooked in $\overline D$ (Lemma~\ref{lem:reverse}), so (a) applies.

(c) Translation is clear. For reflection let $g(x)=f(-x)$ and $x<y$. If $g(y)\ge g(x)+3$, then $f(-y)\ge f(-x)+3$ with $-y<-x$,
so (ii) for $f$ gives $-y<v<v'<-x$ with $f(-x)\le f(v)\le f(-x)+1$ and $f(-y)-2\le f(v')\le f(-y)-1$. Then $u=-v'$, $u'=-v$
satisfy $x<u<u'<y$, $g(y)-2\le g(u)\le g(y)-1$ and $g(x)\le g(u')\le g(x)+1$. The other case is the same with (i) for $f$.
\end{proof}

\formal{\lean{HPC.Fold.Z77.crkOn_of_isCrookedIvF}, \lean{HPC.Fold.Z77.CrkOn.reverse},
\lean{HPC.Fold.B79.minIdx_reverse}.}

\section{Turning levels and folds}\label{sec:fold-folds}

\begin{definition}[Folds]\label{def:fold}
Let $g\colon[0,m]\to\Z$ have unit steps (an \emph{arm}). A \emph{fold} of $g$ onto a function $c\colon\N\to\Z$ (the \emph{walk}) is a
function $\beta\colon[0,m]\to\N$ with $\beta(0)=0$, unit steps, and $\abs{g(j)-c(\beta(j))}\le1$ for all $j$.
\end{definition}

The walks we use are determined by a list of \emph{turning levels} $T_0,T_1,\dots,T_K$ ($K\ge1$) with
$\abs{T_{i+1}-T_i}\ge3$. The \emph{leg} $i$ goes from $T_i$ to $T_{i+1}$. Its \emph{envelope} is the interval
$\env_i=[\min(T_i,T_{i+1})-1,\max(T_i,T_{i+1})+1]$, and its \emph{arrival level} is $\arr_i=T_{i+1}-1$ if $T_i<T_{i+1}$ and
$\arr_i=T_{i+1}+1$ otherwise: the level one short of the end of the leg.

\begin{definition}[Fitting]\label{def:fit}
An arm $g$ \emph{fits the turning levels $T_0,\dots,T_K$ from time $\tau$} if $K\ge1$, $\abs{T_1-T_0}\ge3$, and either
\begin{enumerate}[label=(F\arabic*)]
\item $g(j)\in\env_0$ for all $j\ge\tau$; or
\item there is $\tau'\ge\tau$ with $g(\tau')=\arr_0$, $g(j)\in\env_0$ for $\tau\le j\le\tau'$, and $g$ fits $T_1,\dots,T_K$ from time $\tau'$.
\end{enumerate}
\end{definition}

So the arm runs along the legs: on each leg it stays within one of the levels between the two turning levels, until it
reaches the level just before the end of the leg, where it passes to the next leg; on the last leg it reaches, it stays
for good.

\formal{\lean{HPC.Fold.Z77.Env}, \lean{HPC.Fold.Z77.arr}, \lean{HPC.Fold.Z77.Fit}.}

\paragraph{The jump walk.} Let $\ell_i=\abs{T_{i+1}-T_i}-2\ge1$ and $O_i=\ell_0+\dots+\ell_{i-1}$. Define $c$ on $[0,O_K]$ by
$c(O_i)=T_i$ for $i\le K$ and, for $0<u<\ell_i$,
\[
c(O_i+u)=T_i+2+u\ \text{ if }T_i<T_{i+1},\qquad c(O_i+u)=T_i-2-u\ \text{ if }T_i>T_{i+1}
\]
(and $c(x)=T_K$ for $x>O_K$). So leg $i$ is run as $T_i,T_i\pm3,T_i\pm4,\dots,T_{i+1}\mp1$, and then the next leg starts at
$T_{i+1}$: the walk skips the two levels next to each turning level. For a level $y$ define the \emph{position on leg $i$}
\[
\pi_i(y)=\begin{cases}\min\bigl(\max(y-T_i-1,0),\,\ell_i\bigr)&\text{if }T_i<T_{i+1},\\
\min\bigl(\max(T_i-1-y,0),\,\ell_i\bigr)&\text{if }T_i>T_{i+1}.\end{cases}
\]

\begin{lemma}\label{lem:legpos}
\begin{enumerate}[label=(\alph*)]
\item $0\le\pi_i(y)\le\ell_i$, and $\abs{\pi_i(y)-\pi_i(y')}\le1$ when $\abs{y-y'}\le1$.
\item If $y\in\env_i$, then $\abs{y-c(O_i+\pi_i(y))}\le1$.
\item $\pi_i(y)=0$ if $\abs{y-T_i}\le1$, and $\pi_i(\arr_i)=\ell_i$.
\end{enumerate}
\end{lemma}

\begin{proof}
Suppose $T_i<T_{i+1}$ (the other case is symmetric); write $a=T_i$, $b=T_{i+1}$, so $\ell_i=b-a-2$ and $\env_i=[a-1,b+1]$.
(a) and (c) are immediate from the formula ($\arr_i=b-1$ gives $\pi_i=\min(b-a-2,\ell_i)=\ell_i$). For (b): if $y\le a+1$,
then $\pi_i(y)=0$, $c(O_i)=a$, and $\abs{y-a}\le1$. If $a+2\le y\le b-2$, then $\pi_i(y)=y-a-1\in(0,\ell_i)$ and
$c(O_i+\pi_i(y))=y+1$. If $y\ge b-1$, then $\pi_i(y)=\ell_i$, $c(O_{i+1})=b$, and $\abs{y-b}\le1$ as $y\le b+1$.
\end{proof}

\formal{\lean{HPC.Fold.Z77.jwalk}, \lean{HPC.Fold.Z77.legPos}, \lean{HPC.Fold.Z77.jwalk_legPos},
\lean{HPC.Fold.Z77.legPos_lip}, \lean{HPC.Fold.Z77.legPos_eq_zero}.}

\begin{lemma}[Folds from fitting]\label{lem:fold-of-fit}
Let the arm $g$ fit $T_0,\dots,T_K$ from time $\tau$, with $\abs{g(\tau)-T_0}\le1$. Then there is $\beta\colon[\tau,m]\to[0,O_K]$ with
$\beta(\tau)=0$, unit steps, and $\abs{g(j)-c(\beta(j))}\le1$ for all $j\in[\tau,m]$, where $c$ is the jump walk of $T_0,\dots,T_K$.
Moreover $\beta(j)=0$ whenever $\abs{g(i)-T_0}\le1$ for all $i\in[\tau,j]$.
\end{lemma}

\begin{proof}
Induction on $K$. In case~(F1), let $\beta(j)=\pi_0(g(j))$; by Lemma~\ref{lem:legpos} it starts at $0$, has unit steps, and
satisfies the estimate. In case~(F2), $\abs{g(\tau')-T_1}=\abs{\arr_0-T_1}=1$, so by induction (applied to $T_1,\dots,T_K$, whose
jump walk is $u\mapsto c(\ell_0+u)$) there is $\beta'$ on $[\tau',m]$. Put $\beta(j)=\pi_0(g(j))$ for $j\le\tau'$ and
$\beta(j)=\ell_0+\beta'(j)$ for $j>\tau'$. At $\tau'$ both formulas give $\ell_0$ (as $\pi_0(\arr_0)=\ell_0$ and $\beta'(\tau')=0$),
so $\beta$ has unit steps; the estimate holds on both parts.

For the last statement: $\arr_0$ differs from $T_0$ by at least $2$, so if $g$ stays within one of $T_0$ on $[\tau,j]$, then in
case~(F2) $\tau'>j$, and in both cases $\beta(j)=\pi_0(g(j))=0$ by Lemma~\ref{lem:legpos}(c).
\end{proof}

\formal{\lean{HPC.Fold.Z77.fold_of_fit}, \lean{HPC.Fold.Z77.fold_of_fit_zero}.}

The point is that the walk depends only on the turning levels. If both arms of a sequence fit the same turning levels, their
folds are folds onto the same walk and can be put together.

\section{One phase of the zigzag}\label{sec:fold-phase}

Fix a sequence $S$ on $[-p,q]$ ($p,q\ge0$) with unit steps and bounds $L\le S\le H$. Its two \emph{arms} are $A_+(j)=S(j)$ for
$0\le j\le q$ and $A_-(j)=S(-j)$ for $0\le j\le p$; we write $m_+=q$, $m_-=p$, and $e$ for one of the two signs. Arm times
are $j\ge0$.

\begin{definition}\label{def:phase-state}
An \emph{up state} consists of integers $b,a$ with $b+2\le a$ (the current \emph{bottom} and the previous \emph{top}) and, for
each arm $e$, either nothing (the arm is \emph{finished}) or an \emph{entry time} $\tau_e\le m_e$ with $\abs{A_e(\tau_e)-b}\le1$.
It is \emph{not stuck} if no unfinished arm $e$ has a time $d\ge\tau_e$ with $A_e(d)=b-2$ and $A_e(j)\le a-1$ for $\tau_e\le
j<d$. \emph{Down states} are defined symmetrically: a top $a$, a previous bottom $b$, entry times with $\abs{A_e(\sigma_e)-a}\le1$;
not stuck means that no unfinished arm rises to $a+2$ while staying $\ge b+1$ from its entry time on.
\end{definition}

\begin{lemma}[Up phase]\label{lem:up-phase}
Let an up state be not stuck. Then there is $t\ge a+1$ such that every unfinished arm $e$ either
\begin{itemize}
\item \emph{finishes}: $b-1\le A_e(j)\le t-2$ for all $j\in[\tau_e,m_e]$; or
\item \emph{passes}: there is $\sigma_e>\tau_e$ with $A_e(\sigma_e)=t-1$ and $b-1\le A_e(j)\le t-2$ for $\tau_e\le j<\sigma_e$.
\end{itemize}
Moreover, either every arm finishes, or $t\le H+1$ and some arm $e$ passes and then drops: there is $d>\sigma_e$ with $A_e(d)=b-2$
and $b-1\le A_e(j)\le t-1$ for $\tau_e\le j<d$.
\end{lemma}

\begin{proof}
For each unfinished arm $e$ define a number $v_e$. If $A_e$ goes below $b-1$ at some time after $\tau_e$, let $d_e$ be the first
such time; then $d_e>\tau_e$ (as $A_e(\tau_e)\ge b-1$), $A_e(d_e)=b-2$ by unit steps, and $A_e\ge b-1$ on $[\tau_e,d_e)$. Let $M_e$
be the maximum of $A_e$ on $[\tau_e,d_e)$ and $v_e=M_e+1$. Since the state is not stuck, $M_e\ge a$, so $v_e\ge a+1$; also $v_e\le
H+1$. If $A_e$ never goes below $b-1$, or $e$ is finished, let $v_e=\max(a+1,H+2)$. Let $t=\min(v_+,v_-)\ge a+1$.

Fix an unfinished arm $e$. \emph{Every drop of $A_e$ below $b-1$ after $\tau_e$ is preceded by a visit to $[t-1,\infty)$}: if
$A_e(j)<b-1$, then $d_e\le j$ exists, and $A_e$ takes the value $M_e=v_e-1\ge t-1$ on $[\tau_e,d_e)$. Now if $A_e$ reaches $t-1$
or more after $\tau_e$, let $\sigma_e$ be the first such time. Before $\sigma_e$, $A_e\le t-2$, and $A_e\ge b-1$ (a drop would be
preceded by a visit to $[t-1,\infty)$ before $\sigma_e$); as $A_e(\tau_e)\le b+1\le t-2$, $\sigma_e>\tau_e$ and $A_e(\sigma_e)=t-1$ by
unit steps. So $e$ passes. If $A_e$ never reaches $t-1$, it never drops below $b-1$ either, and it finishes.

If $t=\max(a+1,H+2)$, then $t-1>H$ and no arm can reach $t-1$, so all arms finish. Otherwise $t=v_e<H+2$ for some arm $e$ which
drops at $d_e$, so $t\le H+1$; its maximum $t-1$ on $[\tau_e,d_e)$ is attained, so $e$ passes at a time $\sigma_e<d_e$, and
$b-1\le A_e\le t-1$ on $[\tau_e,d_e)$.
\end{proof}

\formal{\lean{HPC.Fold.Z77.up_step} (with \lean{HPC.Fold.Z77.transfer} and \lean{HPC.Fold.Z77.exists_armVal}).}

\begin{lemma}[Down phase]\label{lem:down-phase}
Let a down state (top $a$, previous bottom $b$) be not stuck. Then there is $t'\le b-1$ such that every unfinished arm either finishes
($t'+2\le A_e\le a+1$ from $\sigma_e$ on) or passes (at a time $\sigma'_e>\sigma_e$ with $A_e(\sigma'_e)=t'+1$, and $t'+2\le A_e\le a+1$
on $[\sigma_e,\sigma'_e)$). Moreover, either every arm finishes, or $t'\ge L-1$ and some arm passes and then rises: there is
$d>\sigma'_e$ with $A_e(d)=a+2$ and $t'+1\le A_e\le a+1$ on $[\sigma_e,d)$.
\end{lemma}

\begin{proof}
Apply Lemma~\ref{lem:up-phase} to the arms $-A_e$, with bottom $-a$, previous top $-b$, upper bound $-L$, and put $t'=-t$.
\end{proof}

\formal{\lean{HPC.Fold.Z77.down_step}.}

Note that these two lemmas do not use crookedness.

\section{Crookedness across the base link}\label{sec:fold-cores}

From now on $S$ is crooked on $[-p,q]$. For an arm $e$ let $\bar e$ be the other arm. Reading $S$ from the end of arm $e$ to the
end of arm $\bar e$ gives the sequence $x\mapsto A_e(-x)$ for $x\le0$, $x\mapsto A_{\bar e}(x)$ for $x\ge0$, which is $S$ or its
reflection, hence crooked (Lemma~\ref{lem:crooked-seq}(c)). In the next two lemmas, the arm called $\mathcal L$ (left) is read at
negative positions and the arm $\mathcal R$ (right) at positive ones; they apply with $\set{\mathcal L,\mathcal R}=\set{A_+,A_-}$ in
either order.

\begin{lemma}[Up core]\label{lem:up-core}
Let $b\le a$. There are no times $\tau\le d_1$ on $\mathcal L$ and $d_2$ on $\mathcal R$ such that $\mathcal L(j)\ge b$ for $j<\tau$,
$\mathcal L(j)\le a-1$ for $\tau\le j<d_1$, $\mathcal L(d_1)=b-2$, $\mathcal R(j)\ge b$ for $j<d_2$, and $\mathcal R(d_2)=a+2$.
\end{lemma}

\begin{proof}
Suppose there are. The positions $x=-d_1<y=d_2$ have values $b-2$ and $a+2$, which differ by at least $4$. By
Definition~\ref{def:crooked-seq}(i) there are $x<u<u'<y$ with value at $u$ in $[a,a+1]$ and value at $u'$ in $[b-2,b-1]$. If $u'\ge0$,
then $u'<d_2$ and $\mathcal R(u')\ge b$, a contradiction. So $u<u'<0$. The arm time $j'=-u'$ has $\mathcal L(j')\le b-1$, so $j'\ge\tau$;
and $j=-u$ satisfies $j'<j<d_1$, so $\mathcal L(j)\le a-1$, contradicting $\mathcal L(j)\ge a$.
\end{proof}

\formal{\lean{HPC.Fold.Z77.up_core}.}

\begin{lemma}[Down core]\label{lem:down-core}
Let $b\le a$. There are no times $\tau\le\sigma\le\rho_1$ on $\mathcal L$ and $\tau'\le d$ on $\mathcal R$ such that
\[
\mathcal L\ge b\text{ on }[0,\tau),\quad \mathcal L\le a-2\text{ on }[\tau,\sigma),\quad \mathcal L\ge b+1\text{ on }[\sigma,\rho_1),\quad
\mathcal L(\rho_1)=a+2,
\]
\[
\mathcal R\ge b\text{ on }[0,\tau'),\quad\mathcal R\le a-1\text{ on }[\tau',d),\quad\mathcal R(d)=b-2 .
\]
\end{lemma}

\begin{proof}
Suppose there are. The positions $x=-\rho_1<y=d$ have values $a+2$ and $b-2$. By Definition~\ref{def:crooked-seq}(ii) there are
$x<u<u'<y$ with value at $u$ in $[b-2,b-1]$ and value at $u'$ in $[a,a+1]$. If $u\ge0$: $\mathcal R(u)\le b-1$ forces $u\ge\tau'$, and
then $u<u'<d$ gives $\mathcal R(u')\le a-1$, a contradiction. So $u<0$, and the arm time $j=-u$ satisfies $\mathcal L(j)\le b-1$; hence
$j\ge\tau$ and $j<\sigma$.

Apply Definition~\ref{def:crooked-seq}(ii) again, to the positions $-\rho_1<u$, whose values $a+2$ and $\mathcal L(j)\le b-1$ differ by at
least $3$: there are $-\rho_1<v<v'<u$ with value at $v$ in $[\mathcal L(j),\mathcal L(j)+1]$, so at most $b$, and value at $v'$ in
$[a,a+1]$. The arm time $-v'$ lies in $(j,\rho_1)$. If $-v'<\sigma$, then $-v'\in[\tau,\sigma)$ and $\mathcal L(-v')\le a-2$, a
contradiction; so $-v'\ge\sigma$, and $-v\in(-v',\rho_1)\subseteq[\sigma,\rho_1)$, whence $\mathcal L(-v)\ge b+1$, a contradiction.
\end{proof}

\formal{\lean{HPC.Fold.Z77.down_core}.}

We now show that the phases of the zigzag never get stuck. An up state has the \emph{history property} if every unfinished arm
satisfies $A_e(j)\ge b$ for $j<\tau_e$.

\begin{lemma}[Up states are not stuck]\label{lem:nofail-up}
Let an up state (bottom $b$, previous top $a$, $b\le a$) have the history property, and suppose some arm $e_0$ and time $d$ satisfy
$A_{e_0}(d)=a+2$, $A_{e_0}(j)\ge b$ for $j<d$, and $\tau_{e_0}\le d$ if $e_0$ is unfinished. Then the state is not stuck.
\end{lemma}

\begin{proof}
Suppose an unfinished arm $e$ has $d_1\ge\tau_e$ with $A_e(d_1)=b-2$ and $A_e\le a-1$ on $[\tau_e,d_1)$. If $e=e_0$: since
$A_e(j)\ge b$ for $j<d$, we have $d<d_1$, so $d\in[\tau_e,d_1)$ and $A_e(d)\le a-1$, contradicting $A_e(d)=a+2$. If $e\ne e_0$:
Lemma~\ref{lem:up-core} with $\mathcal L=A_e$, $\tau=\tau_e$ and $\mathcal R=A_{e_0}$, $d_2=d$ gives a contradiction.
\end{proof}

\formal{\lean{HPC.Fold.Z77.nofail_up}.}

\begin{lemma}[Down states after an up phase are not stuck]\label{lem:nofail-down}
Let an up state (bottom $b$, previous top) have the history property, and let the up phase of Lemma~\ref{lem:up-phase} give the top
$t$, with not every arm finishing. Then the down state with top $t$, previous bottom $b$, and entry times the passing times $\sigma_e$,
is not stuck.
\end{lemma}

\begin{proof}
Let $e_0$ be the arm of Lemma~\ref{lem:up-phase} which passes and then drops: $\tau_{e_0}<\sigma_{e_0}<d$, $A_{e_0}(d)=b-2$, $A_{e_0}\le
t-1$ on $[\tau_{e_0},d)$. Suppose an arm $e$ (which passed) rises to $t+2$ at a time $\rho_1\ge\sigma_e$ while $A_e\ge b+1$ on
$[\sigma_e,\rho_1)$. If $e=e_0$: $\rho_1\ge d$ since $A_{e}\le t-1$ on $[\sigma_e,d)$; but then $d\in[\sigma_e,\rho_1)$ and $A_e(d)=b-2<b+1$.
If $e\ne e_0$: Lemma~\ref{lem:down-core} applies with $\mathcal L=A_e$ (times $\tau_e\le\sigma_e\le\rho_1$; the bound $A_e\le t-2$ on
$[\tau_e,\sigma_e)$ is part of ``passes'', and $A_e\ge b$ before $\tau_e$ is the history property) and $\mathcal R=A_{e_0}$ (times
$\tau_{e_0}\le d$), with $a=t$.
\end{proof}

\formal{\lean{HPC.Fold.Z77.nofail_down}.}

\begin{lemma}[First direction]\label{lem:first-direction}
Let $P=S(0)$. At least one of the following holds:
\begin{enumerate}
\item[$(\mathrm{OK}^-)$] there is no position $x$ with $S(x)=P-3$ and $S\le P$ at all positions between $0$ and $x$;
\item[$(\mathrm{OK}^+)$] there is no position $x$ with $S(x)=P+3$ and $S\ge P$ at all positions between $0$ and $x$.
\end{enumerate}
\end{lemma}

\begin{proof}
Suppose $x_1$ violates the first and $x_2$ the second. If they lie on the same arm, the one nearer to $0$ lies between $0$ and the
other, which is impossible ($S(x_1)=P-3<P$ and $S(x_2)=P+3>P$). Otherwise apply Lemma~\ref{lem:up-core} with $\mathcal L$ the arm of
$x_1$, $\mathcal R$ the arm of $x_2$, $\tau=0$, $a=P+1$, $b=P-1$.
\end{proof}

\formal{\lean{HPC.Fold.Z77.first_direction}; the two conditions are \lean{HPC.Fold.Z77.OKBelow} and \lean{HPC.Fold.Z77.OKAbove}.}

\section{The zigzag theorem}\label{sec:fold-zigzag}

\begin{proposition}[Main induction]\label{prop:main-up}
Let an up state with bottom $b$ and previous top $a$ satisfy $a\le H+1$, $b\ge L-1$, the history property, and be not stuck. Then
there are turning levels $b=T_0,T_1,\dots,T_K$ which every unfinished arm $e$ fits from its entry time $\tau_e$.
\end{proposition}

\begin{proof}
Induction on $(H+1-a)+(b-L+1)\ge0$. Run the up phase (Lemma~\ref{lem:up-phase}): it gives $t\ge a+1\ge b+3$.

\emph{Leg up.} Suppose the passing arms fit some list $t=T'_0,\dots,T'_{K'}$ from their passing times (or that there are no passing
arms and the list is just $t$). Then every unfinished arm fits $b,t,T'_1,\dots$ from $\tau_e$: the envelope of the leg $b\to t$ is
$[b-1,t+1]$ and its arrival level is $t-1$; a finishing arm stays in $[b-1,t-2]$ for good, which is~(F1); a passing arm stays in
$[b-1,t-1]$ on $[\tau_e,\sigma_e]$ and has $A_e(\sigma_e)=t-1$, which is~(F2).

If every arm finishes, the list $b,t$ works. Otherwise, by Lemma~\ref{lem:nofail-down}, the down state with top $t$ and previous
bottom $b$ is not stuck. Run the down phase (Lemma~\ref{lem:down-phase}); it gives $t'\le b-1\le t-4$. The envelope of the leg
$t\to t'$ is $[t'-1,t+1]$ and its arrival level is $t'+1$. An arm finishing in the down phase stays in $[t'+2,t+1]$ for good,
which is~(F1) for this leg; an arm passing it stays in $[t'+1,t+1]$ up to its passing time $\sigma'_e$ and reaches $t'+1$ there. So
if every arm finishes the down phase, the turning levels $b,t,t'$ work; and in general it suffices to find a list starting with $t'$
which the arms passing the down phase fit from their passing times $\sigma'_e$, and to prepend $b,t$.

Otherwise, consider the up state with bottom $t'$, previous top $t$ and entry times $\sigma'_e$. It satisfies $|A_e(\sigma'_e)-t'|\le1$,
$t'+2\le t$, $t\le H+1$ and $t'\ge L-1$ (both from the ``moreover'' parts of the phase lemmas). It has the history property: before
$\sigma'_e$, an arm was $\ge b$ before $\tau_e$ (history), $\ge b-1$ on $[\tau_e,\sigma_e)$, and $\ge t'+2$ on $[\sigma_e,\sigma'_e)$;
all these are $\ge t'$. It is not stuck by Lemma~\ref{lem:nofail-up}, applied with the arm $e_0$ which passes the down phase and then
rises to $t+2$ at time $d$: before $d$ it is $\ge t'$ (by the same history count, and $\ge t'+1$ on $[\sigma_{e_0},d)$), and
$\sigma'_{e_0}\le d$. The measure has decreased by at least $2$, since $t\ge a+1$ and $t'\le b-1$. By induction there is a list
$t',\dots$ which the arms fit from $\sigma'_e$.
\end{proof}

\formal{\lean{HPC.Fold.Z77.main_up}.}

\begin{theorem}[Zigzag theorem]\label{thm:zigzag}
Let $S$ be crooked on $[-p,q]$ with $L\le S\le H$, and let $\epsilon\in\set{+1,-1}$ be such that $(\mathrm{OK}^\epsilon)$ of
Lemma~\ref{lem:first-direction} holds. Then there are a walk $c$ with $c(0)=S(0)+\epsilon$ and a function $\beta\colon[-p,q]\to\N$ with
unit steps and $\beta(0)=0$ such that $\abs{S(x)-c(\beta(x))}\le1$ for all $x$. Moreover $\beta(x)=0$ whenever $\abs{S(y)-c(0)}\le1$ for
all $y$ between $0$ and $x$.
\end{theorem}

\begin{proof}
Let $P=S(0)$; both arms start at $P$. It suffices to find turning levels $T_0=P+\epsilon,T_1,\dots$ which both arms fit from time
$0$: Lemma~\ref{lem:fold-of-fit} then gives folds $\beta_+$, $\beta_-$ of the two arms onto the same jump walk $c$, and
$\beta(x)=\beta_+(x)$ for $x\ge0$, $\beta(x)=\beta_-(-x)$ for $x\le0$ is as required.

If $\epsilon=-1$: the up state with bottom $P-1$, previous top $P+1$, and both entry times $0$, has the history property (vacuously),
$P+1\le H+1$, $P-1\ge L-1$, and it is not stuck by $(\mathrm{OK}^-)$ (an arm dropping to $(P-1)-2=P-3$ while staying $\le P$).
Proposition~\ref{prop:main-up} gives the turning levels, starting at $P-1$.

If $\epsilon=+1$: the down state with top $P+1$ and previous bottom $P-1$ is not stuck by $(\mathrm{OK}^+)$. Run the down phase; it
gives $t'\le P-2$. As in the proof of Proposition~\ref{prop:main-up}, either all arms finish (and the turning levels $P+1,t'$ work) or the
up state with bottom $t'$, previous top $P+1$, and the passing times as entry times, satisfies the hypotheses of
Proposition~\ref{prop:main-up} (history: before passing, the arms are $\ge t'+2$; not stuck by Lemma~\ref{lem:nofail-up}). Prepend $P+1$
to the turning levels it gives.
\end{proof}

\formal{\lean{HPC.Fold.Z77.greedy_fit_below}, \lean{HPC.Fold.Z77.greedy_fit_above}, and \lean{HPC.Fold.Z77.greedy_fit}
(which uses Lemma~\ref{lem:first-direction} to choose $\epsilon$).}

\begin{corollary}[Bing's Theorem~3, three-link form]\label{cor:theorem3}
Let $E$ be crooked in $D$ and let $i$ be an index of $E$. Then there is a block map for $E$ with small blocks in which $E_i$ lies in the
first block.
\end{corollary}

\begin{proof}
Apply Theorem~\ref{thm:zigzag} to $S(x)=\lo(i+x)$ on $[-i,n-1-i]$ (crooked by Lemma~\ref{lem:crooked-seq}, with $0\le S\le k-1$), with
$\epsilon$ from Lemma~\ref{lem:first-direction}. The function $t\mapsto\beta(t-i)$ is a schedule (Definition~\ref{def:schedule}) with value
$0$ at $i$, hence a block map (Lemma~\ref{lem:schedule}). If $\beta(t-i)=u$ then $\abs{\lo(t)-c(u)}\le1$, so $E_t\subseteq
D_{\lo(t)}\subseteq D[c(u)-1,c(u)+1]$.
\end{proof}

\formal{\lean{HPC.Fold.theorem3}, via \lean{HPC.Fold.Z77.s3_zig} and \lean{HPC.exists_blocks_of_walk}.}

\section{Folds with a fake end link}\label{sec:fold-fake}

To prescribe the last block as well as the first, we fold only an initial segment $[0,a]$ of $E$, where $E_a$ meets the first
link $D_0$, and we need to know that the walk comes back down to level $\le1$ after the end of the segment. This is arranged by
appending a \emph{fake link} at level $0$.

In this section $f\colon[0,n-1]\to\Z$ is crooked (Definition~\ref{def:crooked-seq}) with unit steps and $0\le f\le H$, and $i\le a$
are indices. For a \emph{fake level} $F$ consider the two \emph{events}
\begin{itemize}
\item $\mathrm{Ev}^+_F(f,i,a)$: either there is $t\le a$ with $f(t)=f(i)+3$ and $f\ge f(i)$ at all indices strictly between $i$ and
$t$; or $F=f(i)+3$ and $f\ge f(i)$ on $[i,a]$;
\item $\mathrm{Ev}^-_F(f,i,a)$: the same with $f(i)-3$ in place of $f(i)+3$ and $f\le f(i)$ in place of $f\ge f(i)$.
\end{itemize}
So $\mathrm{Ev}^{\pm}_F(f,i,a)$ says that the sequence $f$ on $[0,a]$, continued by one more value $F$ at $a+1$, violates the condition
$(\mathrm{OK}^\pm)$ of Lemma~\ref{lem:first-direction} at the base $i$.

\formal{\lean{HPC.Fold.B79.EvUp} and \lean{HPC.Fold.B79.EvDown}; reflecting $f$ exchanges them
(\lean{HPC.Fold.B79.evUp_reflect}).}

\begin{lemma}[Crooked extensions]\label{lem:crooked-ext}
\begin{enumerate}[label=(\alph*)]
\item Suppose $f(a)\le1$ and the \emph{low meeting property} holds: for every $x<a$ with $f(x)\ge3$ there are $x<y<y'<a$ with
$f(y)\le1$ and $f(x)-2\le f(y')\le f(x)-1$. Then the sequence $f^\flat$ on $[0,a+1]$ equal to $f$ on $[0,a]$ with $f^\flat(a+1)=0$ is
crooked, has unit steps, and satisfies $0\le f^\flat\le H$.
\item Suppose $f(a)\ge H-1$ and the \emph{high meeting property} holds: for every $x<a$ with $f(x)\le H-3$ there are $x<y<y'<a$ with
$H-2\le f(y)\le H-1$ and $f(x)\le f(y')\le f(x)+1$. Then the sequence $f^\sharp$ equal to $f$ on $[0,a]$ with $f^\sharp(a+1)=H$ is
crooked, has unit steps, and satisfies $0\le f^\sharp\le H$.
\end{enumerate}
\end{lemma}

\begin{proof}
(a) The step from $a$ to $a+1$ is $f(a)-0\in\set{0,1}$, and the bounds are clear. Pairs $x<y\le a$ are covered by crookedness of
$f$. Let $y=a+1$ and $x\le a$. Clause~(i) of Definition~\ref{def:crooked-seq} needs $0\ge f(x)+3$, impossible as $f\ge0$. For
clause~(ii) let $0\le f(x)-3$; then $x\ne a$ (as $f(a)\le1$), so $x<a$, and the low meeting property gives $x<u<u'<a$ with $0\le
f(u)\le1=f^\flat(a+1)+1$ and $f(x)-2\le f(u')\le f(x)-1$, which is clause~(ii) for the pair $(x,a+1)$.

(b) The same with the roles of the two clauses exchanged: clause~(ii) for $(x,a+1)$ needs $H\le f(x)-3$, impossible; clause~(i)
needs $H\ge f(x)+3$, so $x<a$, and the high meeting property supplies the two indices.
\end{proof}

\formal{\lean{HPC.Fold.B79.crkOn_lowExt}, \lean{HPC.Fold.B79.crkOn_highExt} (on sequences re-indexed so that the base is at $0$).}

\begin{lemma}[Fold with a fake low end]\label{lem:fake-low}
Assume the hypotheses of Lemma~\ref{lem:crooked-ext}(a), and let $\epsilon\in\set{+1,-1}$ be such that $\mathrm{Ev}^\epsilon_0(f,i,a)$
fails. Then there are a walk $c\colon\N\to\Z$, a function $\beta\colon[0,a]\to\N$ and a number $U\in\N$ such that
\begin{enumerate}[label=(\alph*)]
\item $c(0)=f(i)+\epsilon$, $\beta(i)=0$, $\beta$ has unit steps, and $\abs{f(t)-c(\beta(t))}\le1$ for $t\le a$;
\item $\abs{U-\beta(a)}\le1$ and $c(U)\le1$;
\item if $t\le a$ is a neighbour of $i$ ($t=i\pm1$) with $\abs{f(t)-c(0)}\le1$, then $\beta(t)=0$;
\item if $a=i$ and $\abs{c(0)}\le1$, then $U=0$.
\end{enumerate}
\end{lemma}

\begin{proof}
Let $S(x)=f^\flat(i+x)$ for $-i\le x\le a-i+1$. By Lemma~\ref{lem:crooked-ext}(a) and Lemma~\ref{lem:crooked-seq}(c), $S$ is crooked
on $[-i,a-i+1]$, with unit steps and $0\le S\le H$; $S(0)=f(i)$. The condition $(\mathrm{OK}^\epsilon)$ of
Lemma~\ref{lem:first-direction} for $S$ is the failure of $\mathrm{Ev}^\epsilon_0(f,i,a)$: a position $x\le a-i$ violating it is an index
$t=i+x\le a$ as in the first alternative of the event, and the position $a-i+1$ (value $0$) violates it exactly when the second
alternative holds, the positions strictly between $0$ and $a-i+1$ being those of $[i,a]$. So Theorem~\ref{thm:zigzag} applies and
gives $c$ and $\tilde\beta$ on $[-i,a-i+1]$. Put $\beta(t)=\tilde\beta(t-i)$ and $U=\tilde\beta(a-i+1)$. Then (a) is the statement of
the theorem; (b) follows from unit steps and $\abs{0-c(U)}=\abs{S(a-i+1)-c(\tilde\beta(a-i+1))}\le1$. For (c), the positions between $0$
and $t-i=\pm1$ are $0$ and $\pm1$, and $\abs{S(0)-c(0)}=1$; so the ``moreover'' part of Theorem~\ref{thm:zigzag} gives
$\tilde\beta(t-i)=0$. For (d), if $a=i$ the positions between $0$ and $1$ are $0$ and $1$, with $S(1)=0$; as $\abs{c(0)}\le1$, again
$U=\tilde\beta(1)=0$.
\end{proof}

\formal{\lean{HPC.Fold.B79.finfold_low}, from \lean{HPC.Fold.B79.fold_low_end_mode} (the version on $\Z$-indexed sequences;
without the choice of $\epsilon$ it is \lean{HPC.Fold.B78.fold_low_end}). With $\epsilon$ chosen by Lemma~\ref{lem:first-direction}
it is \lean{HPC.Fold.B79.finfold_low_any}.}

\begin{lemma}[Fold with a fake high end]\label{lem:fake-high}
Assume the hypotheses of Lemma~\ref{lem:crooked-ext}(b), and let $\epsilon$ be such that $\mathrm{Ev}^\epsilon_H(f,i,a)$ fails. Then
there are $c$, $\beta$, $U$ with properties (a), (c) of Lemma~\ref{lem:fake-low} and $\abs{U-\beta(a)}\le1$, $c(U)\ge H-1$.
\end{lemma}

\begin{proof}
The same, with $f^\sharp$ and Lemma~\ref{lem:crooked-ext}(b): now $\abs{H-c(U)}\le1$.
\end{proof}

\formal{\lean{HPC.Fold.B79.finfold_high}, from \lean{HPC.Fold.B79.fold_high_end_mode}.}

\paragraph{The meeting properties.} They come from crookedness when $E_a$ meets $D_0$.

\begin{lemma}\label{lem:meet-low}
Let $E$ be crooked in $D$ (of length $k\ge2$), let $E_a$ meet $D_0$, let $t<a$, and let $E_t\subseteq D_j$ with $j\ge3$. Then there are
$t<y<y'<a$ with $E_y\subseteq D_1$ and $E_{y'}\subseteq D_{j-1}$. Consequently:
\begin{enumerate}[label=(\alph*)]
\item $\lo$ has the low meeting property at $a$, and $\lo(a)\le1$;
\item $\bar\lo=k-1-\hi$ has the high meeting property at $a$ with $H=k-1$, and $\bar\lo(a)\ge k-2$;
\item for every $t<a$ with $\lo(t)\ge3$ there are $t<y<y'<a$ with $\lo(y)\le1$ and $\hi(y')\ge\lo(t)-1$.
\end{enumerate}
\end{lemma}

\begin{proof}
$E_a$ meets $D_0$, $E_t$ meets $D_j$, and $j-0\ge3$; crookedness gives indices $r,s$ with $\mathrm{Btw}(a,r,s,t)$, $E_r\subseteq D_{j-1}$
and $E_s\subseteq D_1$. Since $t<a$, $\mathrm{Btw}(a,r,s,t)$ means $t<s<r<a$; take $y=s$, $y'=r$.

(a) Apply this with $j=\lo(t)$. Then $\lo(y)\le1$, and $D_{\lo(t)-1}\in\set{D_{\lo(y')},D_{\hi(y')}}$, so $\lo(t)-2\le\lo(y')\le\lo(t)-1$
(Lemma~\ref{lem:readings}(a)). Also $\lo(a)\le1$ and $\hi(a)\le1$ by Lemma~\ref{lem:readings}(c).
(b) Apply it with $j=\hi(t)$, which is $\ge3$ when $\bar\lo(t)\le k-4$. Then $\hi(y)\in\set{1,2}$, i.e.\ $k-3\le\bar\lo(y)\le k-2$, and
$\hi(y')\in\set{\hi(t)-1,\hi(t)}$, i.e.\ $\bar\lo(t)\le\bar\lo(y')\le\bar\lo(t)+1$. And $\hi(a)\le1$ gives $\bar\lo(a)\ge k-2$.
(c) Apply it with $j=\lo(t)$: $\hi(y')\ge\lo(t)-1$ as $E_{y'}\subseteq D_{\lo(t)-1}$.
\end{proof}

\formal{\lean{HPC.Fold.B79.meet_low_gen}, \lean{HPC.Fold.B79.maxIdx_near_zero}; the two sequences and their properties are
collected in \lean{HPC.Fold.B79.seq_facts}.}

\paragraph{Exclusion.} The prefix fold must send \emph{both} links $r,r+1$ to time $0$. When $\lo(r)\ne\lo(r+1)$ this can fail for
the fold read on $\lo$ in the only direction which would work; then we read on $\hi$ instead. The next lemma says that one of the two
readings always works.

\begin{lemma}[Exclusion]\label{lem:exclusion}
Let $E$ be crooked in $D$, let $E_a$ meet $D_0$, and let $\ell,h$ be adjacent indices (that is, $\abs{\ell-h}=1$) with $\ell\le a$ and
\[
x:=\lo(\ell),\qquad \lo(h)=\hi(h)=x+1 .
\]
Then $\mathrm{Ev}^+_0(\lo,\ell,a)$ and $\mathrm{Ev}^-_0(\hi,h,a)$ do not both hold.
\end{lemma}

\begin{proof}
Suppose both hold. The second alternative of $\mathrm{Ev}^+_0(\lo,\ell,a)$ would need $x+3=0$, impossible; so there is $t\le a$ with
$\lo(t)=x+3$ and $\lo\ge x$ strictly between $\ell$ and $t$. Note $t\ne\ell,h$ (their $\lo$ is $x$, $x+1$), and $\hi(t)\ge x+3$. We
use repeatedly that an index strictly between $\ell$ and $t$ has $\lo\ge x$, and the following consequence of adjacency of $\ell,h$:
\emph{an index $w\ne\ell,h$ lies on the same side of $\ell$ as of $h$}.

\emph{Case 1: the first alternative of $\mathrm{Ev}^-_0(\hi,h,a)$}: there is $u\le a$ with $\hi(u)=x-2$ and $\hi\le x+1$ strictly between
$h$ and $u$. Then $u\ne\ell,h$ (their $\hi$ is $\ge x$) and $\lo(u)\le x-2$.
\begin{itemize}
\item $t,u$ both below the pair ($t<\ell$, $u<h$). If $t<u$, then $u$ is strictly between $t$ and $\ell$, so $\lo(u)\ge x$, a contradiction.
If $u<t$, then $t$ is strictly between $u$ and $h$, so $\hi(t)\le x+1$, a contradiction. The case of both above is the same.
\item $t$ below, $u$ above ($t<\ell$, $u>h$). $E_t$ meets $D_{x+3}$ and $E_u$ meets $D_{x-2}$; crookedness gives $t<w<w'<u$ with
$E_w\subseteq D_{x-1}$ and $E_{w'}\subseteq D_{x+2}$. As $\lo(w)\le x-1$, $w$ is not strictly between $t$ and $\ell$ and $w\ne\ell$; so
$w>\ell$, and then $w'>\ell$. As $\hi(w')\ge x+2$, $w'\ne h$, so $w'$ is strictly between $h$ and $u$, and $\hi(w')\le x+1$, a contradiction.
\item $t$ above, $u$ below. Crookedness for $u<t$ ($E_u$ meets $D_{x-2}$, $E_t$ meets $D_{x+3}$) gives $u<w<w'<t$ with $E_w\subseteq
D_{x+2}$ and $E_{w'}\subseteq D_{x-1}$. As $\hi(w)\ge x+2$, $w$ is not strictly between $u$ and $h$ and $w\ne h$, so $w>h$ and $w'>h$; as
$\lo(w')\le x-1$, $w'\ne\ell$, so $w'$ is strictly between $\ell$ and $t$ and $\lo(w')\ge x$, a contradiction.
\end{itemize}

\emph{Case 2: the second alternative}: $x-2=0$ and $\hi\le x+1=3$ on $[h,a]$. If $t>\ell$, then $t\ge h$ (as $t\ne h$) and $t\le a$, so
$\hi(t)\le3<5=\lo(t)$, impossible. If $t<\ell$: $t<a$ and $\lo(t)=5\ge3$, so Lemma~\ref{lem:meet-low}(c) gives $t<y<y'<a$ with $\lo(y)\le1$
and $\hi(y')\ge4$. As $\lo(y)<x$, $y$ is not strictly between $t$ and $\ell$, and $y\ne\ell$; so $y>\ell$ and $y'\ge h$. Then $y'\in[h,a]$ and
$\hi(y')\le3$, a contradiction.
\end{proof}

\formal{\lean{HPC.Fold.B79.pair_exclusion}.}

\begin{proposition}[Prefix fold]\label{prop:prefix-fold}
Let $k\ge4$, let $E$ be crooked in $D$, let $r'=r+1$, and let $a\ge r$ be an index such that $E_a$ meets $D_0$. Then there are a
\emph{reading} $\lambda\colon[0,n-1]\to\Z$ with $\lo\le\lambda\le\hi$, a walk $c$, a function $\beta\colon[0,a]\to\N$ with unit steps
and a number $U$ such that
\begin{enumerate}[label=(\alph*)]
\item $\beta(r)=0$, and either $r'\le a$ and $\beta(r')=0$, or $a=r$ and $U=0$;
\item $\abs{\lambda(t)-c(\beta(t))}\le1$ for $t\le a$;
\item $\abs{U-\beta(a)}\le1$ and $c(U)\le1$.
\end{enumerate}
\end{proposition}

\begin{proof}
By Lemma~\ref{lem:meet-low}(a), Lemma~\ref{lem:fake-low} applies to $f=\lo$ (crooked by Lemma~\ref{lem:crooked-seq}, $0\le\lo\le k$)
with any base $i\le a$.

\emph{Case $a=r$.} Take $\lambda=\lo$, base $r$, and $\epsilon=-1$: the event $\mathrm{Ev}^-_0(\lo,r,a)$ fails, because $\lo(r)=\lo(a)\le1$
and neither a value $\lo(t)=\lo(r)-3<0$ nor $0=\lo(r)-3$ is possible. Then $c(0)=\lo(r)-1\in\set{-1,0}$, so $U=0$ by
Lemma~\ref{lem:fake-low}(d).

From now on $a>r$, so $r'\le a$. \emph{Case $\lo(r')=\lo(r)$.} Take $\lambda=\lo$, base $r$, and $\epsilon$ given by
Lemma~\ref{lem:first-direction} for the extended sequence (one of the two events fails). Then $\abs{\lo(r')-c(0)}=1$, so $\beta(r')=0$ by
Lemma~\ref{lem:fake-low}(c).

\emph{Case $\lo(r')\ne\lo(r)$.} Then they differ by one; let $\ell\in\set{r,r'}$ be the one with the smaller value $x=\lo(\ell)$ and $h$
the other, so $\lo(h)=x+1$. Also $\hi(h)=x+1$: $\hi(h)\ge\lo(h)$, and $\hi(h)\le\lo(\ell)+1$ by Lemma~\ref{lem:readings}(b). It suffices
to find $\lambda,c,\beta,U$ with (b), (c) and $\beta(\ell)=\beta(h)=0$.
\begin{itemize}
\item If $\mathrm{Ev}^+_0(\lo,\ell,a)$ fails, take $\lambda=\lo$, base $\ell$, $\epsilon=+1$. Then $c(0)=x+1=\lo(h)$, so $\beta(h)=0$ by
Lemma~\ref{lem:fake-low}(c).
\item Otherwise $\mathrm{Ev}^-_0(\hi,h,a)$ fails by Lemma~\ref{lem:exclusion}. Reflecting, $\mathrm{Ev}^+_{k-1}(\bar\lo,h,a)$ fails for
$\bar\lo=k-1-\hi$. By Lemma~\ref{lem:meet-low}(b), Lemma~\ref{lem:fake-high} applies to $\bar\lo$ (crooked by
Lemma~\ref{lem:crooked-seq}(b)) with $H=k-1$, base $h$ and $\epsilon=+1$, giving $\bar c,\beta,U$ with $\bar c(0)=\bar\lo(h)+1$ and
$\bar c(U)\ge k-2$. Take $\lambda=\hi$ and $c=k-1-\bar c$. Then (b) holds, $c(U)\le1$, and $c(0)=\hi(h)-1=x$; since
$\hi(\ell)\in\set{x,x+1}$, $\abs{\bar\lo(\ell)-\bar c(0)}=\abs{\hi(\ell)-x}\le1$, so $\beta(\ell)=0$ by Lemma~\ref{lem:fake-low}(c).\qedhere
\end{itemize}
\end{proof}

\formal{\lean{HPC.Fold.B79.pair_prefix_fold}; the conclusion is \lean{HPC.Fold.B79.PrefixFold}.}

\begin{corollary}[Suffix fold]\label{cor:suffix-fold}
Let $k\ge4$, let $E$ be crooked in $D$, let $s=s'+1$, and let $b\le s$ be an index such that $E_b$ meets $D_{k-1}$. Then there are a
reading $\mu$ with $\lo\le\mu\le\hi$, a walk $c'$, a function $\beta'\colon[b,n-1]\to\N$ with unit steps, and $V$, such that
$\beta'(s)=0$; either $b\le s'$ and $\beta'(s')=0$, or $b=s$ and $V=0$; $\abs{(k-1-\mu(t))-c'(\beta'(t))}\le1$ for $t\ge b$; and
$\abs{V-\beta'(b)}\le1$, $c'(V)\le1$.
\end{corollary}

\begin{proof}
Apply Proposition~\ref{prop:prefix-fold} to the reversed chains: $\overline E$ is crooked in $\overline D$ (Lemma~\ref{lem:reverse},
twice), the pair $s,s'$ becomes the pair $n-1-s$, $n-1-s'=(n-1-s)+1$, and $\overline E_{n-1-b}=E_b$ meets $\overline D_0=D_{k-1}$. The
least and greatest indices of $\overline E_{n-1-t}$ in $\overline D$ are $k-1-\hi(t)$ and $k-1-\lo(t)$, so a reading $\bar\lambda$ for the
reversed chains gives the reading $\mu(t)=k-1-\bar\lambda(n-1-t)$ of $E$ in $D$.
\end{proof}

\formal{\lean{HPC.Fold.B79.pair_suffix_fold}.}

\section{The splice}\label{sec:fold-splice}

We now combine a prefix fold on $[0,a]$ and a suffix fold on $[b,n-1]$ with the middle part of $E$, which is grouped by its levels.

\begin{proposition}[Splice]\label{prop:splice}
Let $k\ge6$, let $E$ be crooked in $D$, and let $r\le a<b\le s$ with $\lo(a)\le1$ and $\hi(b)\ge k-2$. Let
$(\lambda,c,\beta_P,U)$ satisfy (b), (c) of Proposition~\ref{prop:prefix-fold} with $\beta_P(r)=0$, and let $(\mu,c',\beta_Q,V)$ satisfy
the conclusion of Corollary~\ref{cor:suffix-fold} except the clause about $s'$. Then there is a schedule $\beta\colon[0,n-1]\to\N$ (unit
steps) such that
\begin{enumerate}[label=(\alph*)]
\item $\beta(r)=0$ and $\beta(t)\le\beta(s)$ for all $t$;
\item every block $\set{t:\beta(t)=v}$ lies in some $D[w,w+2]$: there is $w$ such that for each $t$ with $\beta(t)=v$, $\lo(t)$ or $\hi(t)$
lies in $[w,w+2]$;
\item if $t\le a$ and $\beta_P(t)=0$, then $\beta(t)=0$; if $U=0$, $t<b$ and $\lo(t)\le2$, then $\beta(t)=0$;
\item if $t\ge b$ and $\beta_Q(t)=0$, then $\beta(t)=\beta(s)$; if $V=0$, $t>a$ and $\lo(t)\ge k-3$, then $\beta(t)=\beta(s)$.
\end{enumerate}
\end{proposition}

Note that ``$\lo(t)$ or $\hi(t)$ in $[w,w+2]$'' implies $E_t\subseteq D[w,w+2]$, since $E_t\subseteq D_{\lo(t)}$ and $E_t\subseteq
D_{\hi(t)}$.

\begin{proof}
Let $G=k-5\ge1$ and $\gamma(y)=\min(\max(y-2,0),G)$; $\gamma$ is non-decreasing and moves by at most one when $y$ does. The \emph{middle
schedule} is $g(t)=\gamma(\lo(t))$ for $t<b$ and $g(t)=\gamma(\hi(t))$ for $t\ge b$. Let $K_P=\set{t\le a:\beta_P(t)<U}$ (the part of the
prefix fold kept) and $K_S=\set{t\ge b:\beta_Q(t)<V}$, and define
\[
\beta(t)=\begin{cases}\beta_P(t)&t\in K_P,\\ U+G+V-\beta_Q(t)&t\in K_S,\\ U+g(t)&\text{otherwise.}\end{cases}
\]
We first note $a+3\le b$: by Lemma~\ref{lem:readings}(b), if $b=a+1$ then $\hi(b)\le\lo(a)+1\le2$, and if $b=a+2$ then $\hi(b)\le\lo(a+1)+1
\le\lo(a)+2\le3$; both contradict $\hi(b)\ge k-2\ge4$. So $K_P$ and $K_S$ are disjoint, and far apart.

Two observations do most of the work. \emph{(P)} If $t\le a$ and $\beta_P(t)=U$, then $\abs{\lambda(t)-c(U)}\le1$ and $c(U)\le1$, so
$\lo(t)\le\lambda(t)\le2$ and $g(t)=\gamma(\lo(t))=0$. \emph{(Q)} If $t\ge b$ and $\beta_Q(t)=V$, then $k-1-\mu(t)\le c'(V)+1\le2$, so
$\hi(t)\ge\mu(t)\ge k-3$ and $g(t)=\gamma(\hi(t))=G$ (as $\hi(t)-2\ge k-5$).

\emph{Unit steps.} Let $t'=t+1$.
\begin{itemize}
\item Both in $K_P$, or both in $K_S$: unit steps of $\beta_P$, $\beta_Q$.
\item $t\in K_P$, $t'\notin K_P$ (so $t'\notin K_S$, as $t'\le a+1<b$). If $t'\le a$, then $\beta_P(t)<U\le\beta_P(t')$ forces $\beta_P(t)=U-1$,
$\beta_P(t')=U$, and by (P) $\beta(t')=U$. If $t'=a+1$, then $t=a$, $\beta_P(a)=U-1$ (from $\abs{U-\beta_P(a)}\le1$), and
$\lo(t')\le\lo(a)+1\le2$, so $\beta(t')=U$. In both cases $\beta(t)=U-1$. Entering $K_P$ ($t\notin K_P$, $t'\in K_P$) is the same.
\item $t\in K_S$, $t'\notin K_S$: then $t'\ge b$, $\beta_Q(t)=V-1$, $\beta_Q(t')=V$, and by (Q) $\beta(t')=U+G=\beta(t)-1$. Entering $K_S$ at
$t'$: if $t\ge b$ this is the same; otherwise $t'=b$, $\beta_Q(b)=V-1$ (from $\abs{V-\beta_Q(b)}\le1$), $\beta(t')=U+G+1$, and
$\lo(t)\ge\hi(b)-1\ge k-3$ by Lemma~\ref{lem:readings}(b), so $\beta(t)=U+G$.
\item Neither in $K_P\cup K_S$: $\beta=U+g$, and $g$ has unit steps since $\gamma$ is $1$-Lipschitz and $\lo$, $\hi$ have unit steps;
at the switch $t=b-1$, $t'=b$, use $\abs{\hi(b)-\lo(b-1)}\le1$.
\end{itemize}

\emph{(a)} If $U>0$, then $r\in K_P$ and $\beta(r)=\beta_P(r)=0$. If $U=0$, then $r\notin K_P\cup K_S$, and by (P) (with $\beta_P(r)=0=U$)
$\beta(r)=g(r)=0$. Clearly $\beta\le U+G+V$ everywhere ($\beta<U$ on $K_P$, $g\le G$). And $\beta(s)=U+G+V$: if $V>0$, $s\in K_S$ and
$\beta_Q(s)=0$; if $V=0$, $s\notin K_S$ and by (Q) $\beta(s)=U+G$.

\emph{(b)} Let $v<U$. Then every $t$ with $\beta(t)=v$ lies in $K_P$ (outside $K_P$, $\beta\ge U$) and has $\beta_P(t)=v$, so
$\lambda(t)\in[c(v)-1,c(v)+1]$; since $\lambda(t)$ is $\lo(t)$ or $\hi(t)$ (Lemma~\ref{lem:readings}(a)), take $w=c(v)-1$. Let $U\le
v\le U+G$ and $j=v-U$; then $\beta(t)=v$ forces $t\notin K_P\cup K_S$ and $g(t)=j$, so the relevant reading $y$ ($\lo(t)$ if $t<b$, $\hi(t)$
otherwise) satisfies $\gamma(y)=j$: if $j=0$, $0\le y\le2$ (take $w=0$); if $0<j<G$, $y=j+2$ (take $w=j$); if $j=G$, $k-3\le y\le k-1$ (take
$w=k-3$). Let $v>U+G$; then $t\in K_S$ and $\beta_Q(t)=U+G+V-v$, so $k-1-\mu(t)$ is within one of $c'(U+G+V-v)$; take $w=k-2-c'(U+G+V-v)$.

\emph{(c)} If $t\le a$ and $\beta_P(t)=0$: if $U>0$, $t\in K_P$ and $\beta(t)=0$; if $U=0$, by (P) $\beta(t)=g(t)=0$. If $U=0$, $t<b$ and
$\lo(t)\le2$: $t\notin K_P\cup K_S$ and $g(t)=\gamma(\lo(t))=0$.

\emph{(d)} If $t\ge b$ and $\beta_Q(t)=0$: if $V>0$, $t\in K_S$ and $\beta(t)=U+G+V$; if $V=0$, by (Q) $\beta(t)=U+G=\beta(s)$. If $V=0$,
$t>a$ and $\lo(t)\ge k-3$: $t\notin K_P\cup K_S$, and $g(t)=G$ since $\hi(t)\ge\lo(t)\ge k-3$.
\end{proof}

\formal{\lean{HPC.Fold.B79.splice_schedule_pair} (the version with $\lambda=\mu=\lo$ and without (c), (d) is
\lean{HPC.Fold.B78.splice_schedule}); the passage from levels to links is \lean{HPC.Fold.B78.link_sub_of_level}.}

\section{Proof of the folding theorem}\label{sec:fold-proof}

\begin{lemma}[From a schedule to the block map]\label{lem:assemble}
Let $\beta$ be a schedule for $E$ with $\beta(r)=\beta(r')=0$, $\beta(s')=\beta(s)\ge\beta(t)$ for all $t$, and such that every block of
$\beta$ lies in some $D[w,w+2]$. Then the conclusion of Theorem~\ref{thm:folding} holds.
\end{lemma}

\begin{proof}
By Lemma~\ref{lem:schedule}, $\beta$ is a block map onto $[q]$ with $q=1+\beta(s)$. Its blocks are those of $\beta$, and $\beta(s)=\beta(s')=q-1$.
\end{proof}

\formal{\lean{HPC.Fold.B79.assemble_pair}.}

\begin{lemma}[Two windows]\label{lem:windows}
Let $A,B$ be sets of indices with $A\cup B=[0,n-1]$, such that every $E_t$ with $t\in A$ lies in $D[w_A,w_A+2]$ and every $E_t$ with $t\in
B$ lies in $D[w_B,w_B+2]$. Let $r'=r+1<s'=s-1$ with $r,r'\in A$, and suppose that either $s',s\in B$, or $s',s\in A$ and some $c$ with
$r'<c<s'$ lies in $B$. Then the conclusion of Theorem~\ref{thm:folding} holds.
\end{lemma}

\begin{proof}
In the first case let $\beta(t)=1$ for $t\in\set{s',s}$ and otherwise $\beta(t)=0$ for $t\in A$, $\beta(t)=1$ for $t\notin A$. In the second
case let, for $t<c$, $\beta(t)=0$ if $t\in A$ and $1$ otherwise; $\beta(c)=1$; and, for $t>c$, $\beta(t)=2$ if $t\in A$ and $1$ otherwise.
In both cases $\beta$ has unit steps (in the second case the only way to jump from $0$ to $2$ would be across $c$, where $\beta=1$),
$\beta(r)=\beta(r')=0$, and $\beta(s)=\beta(s')$ is the maximum. Block $1$ consists of indices outside $A$, hence in $B$, together with
$s',s$, resp.\ $c$, which are in $B$; the other blocks consist of indices in $A$. Apply Lemma~\ref{lem:assemble}.
\end{proof}

\formal{\lean{HPC.Fold.B79.pair_windows}.}

\begin{proof}[Proof of Theorem~\ref{thm:folding}]
Let $h_l$ and $h_h$ be the given indices: $E_{h_l}$ meets $D_0$ and $E_{h_h}$ meets $D_{k-1}$, so $\lo(h_l)\le1$ and $\hi(h_h)\ge k-2$
(Lemma~\ref{lem:readings}(c)). Let $A=\set{t:\lo(t)\le2}$ and $B=\set{t:\hi(t)\ge k-3}$; then $E_t\subseteq D[0,2]$ for $t\in A$ and
$E_t\subseteq D[k-3,k-1]$ for $t\in B$.

\emph{$k=3$.} Every link lies in $D[0,2]=\bigcup D$; take the single block, $\beta\equiv0$, $q=1$.

\emph{$k\ge4$: then $r'<s'$.} Otherwise $s=r+2$, and $h_l,h_h\in[r,r+2]$. Crookedness applied to $E_{h_l}$ (meeting $D_0$) and $E_{h_h}$ (meeting
$D_{k-1}$, $k-1\ge3$) gives two indices strictly between $h_l$ and $h_h$, which is impossible within three consecutive indices.

\emph{$k\ge6$.} $h_l\ne h_h$, since $\hi(h_l)\le\lo(h_l)+1\le2<k-2$. If $h_l<h_h$, let $a=h_l$, $b=h_h$, so $r\le a<b\le s$. Take the prefix
fold of Proposition~\ref{prop:prefix-fold} for $(r,r',a)$ and the suffix fold of Corollary~\ref{cor:suffix-fold} for $(s',s,b)$, and splice
them (Proposition~\ref{prop:splice}). The schedule $\beta$ satisfies $\beta(r)=0$ and $\beta(t)\le\beta(s)$, and its blocks are small.
Moreover $\beta(r')=0$: either $r'\le a$ and $\beta_P(r')=0$, and Proposition~\ref{prop:splice}(c) applies; or $a=r$ and $U=0$, and then
$r'=a+1<b$ and $\lo(r')\le\lo(a)+1\le2$, so again (c) applies. Symmetrically $\beta(s')=\beta(s)$ by (d). Lemma~\ref{lem:assemble} finishes
this case. If $h_h<h_l$, apply this to $E$ and the reversed chain $\overline D$ (in which $E_{h_h}$ meets $\overline D_0$ and $E_{h_l}$ meets
$\overline D_{k-1}$; $E$ is crooked in $\overline D$ by Lemma~\ref{lem:reverse}); a block inside $\overline D[j,j+2]$ lies inside
$D[k-3-j,k-1-j]$.

\emph{$k\in\set{4,5}$.} Every index lies in $A\cup B$: if $\lo(t)\ge3$ then $\hi(t)\ge3\ge k-3$. Also, an index $t$ not in $B$ has
$\hi(t)\le k-4\le1$, so its neighbours have $\lo\le2$: \emph{if one of two adjacent indices is not in $B$, both are in $A$}; similarly, an
index not in $A$ has $\lo\ge3$, so its neighbours have $\hi\ge2\ge k-3$: \emph{if one of two adjacent indices is not in $A$, both are in
$B$}. Hence each of the pairs $\set{r,r'}$ and $\set{s',s}$ lies in $A$ or in $B$.
\begin{itemize}
\item If $r,r'\in A$ and $s',s\in B$, apply Lemma~\ref{lem:windows} (first case). If $r,r'\in B$ and $s',s\in A$, apply it with $A$ and $B$
exchanged.
\item Otherwise, suppose $r,r'\in A$. Then $s',s\in A$ (since not both are in $B$), and $h_h$ lies strictly between $r'$ and $s'$: if
$h_h\in\set{r,r'}$, then $\hi\ge k-3$ on both $r,r'$ (as $\hi(h_h)\ge k-2$), so $r,r'\in B$ and we would be in the previous case; if
$h_h\in\set{s',s}$, similarly $s',s\in B$. Since $h_h\in B$, Lemma~\ref{lem:windows} (second case) applies with $c=h_h$.
\item Otherwise $r,r'$ are not both in $A$, so both are in $B$, and then $s',s$ are not both in $A$, so both are in $B$. Now $h_l\in A$
lies strictly between $r'$ and $s'$: if $h_l\in\set{r,r'}$, then $\lo\le2$ on both $r,r'$, i.e.\ $r,r'\in A$; if $h_l\in\set{s',s}$, then
similarly $s',s\in A$, a contradiction. Apply Lemma~\ref{lem:windows} (second case) with $A$ and $B$ exchanged and $c=h_l$.\qedhere
\end{itemize}
\end{proof}

\formal{\lean{HPC.Fold.theorem4_pair}; the main case is \lean{HPC.Fold.B79.theorem4_pair_core}, and the reflection uses
\lean{HPC.Fold.B78.unionRange_reverse}.}


\section*{Exercises}
\addcontentsline{toc}{section}{Exercises}

\begin{exercise}
Show that no sequence beginning with the values $0,1,2,3$ is crooked. Show that every sequence with unit steps whose values lie in an
interval of length at most $2$ is crooked.
\end{exercise}

\begin{exercise}
Compute the jump walk of the turning levels $0,4,1,5$ and the leg positions $\pi_0(y)$ for $-1\le y\le5$. Check
Lemma~\ref{lem:legpos} directly in this case.
\end{exercise}

\begin{exercise}
Let $S(x)=x$ on $[-3,3]$. Show that both conditions $(\mathrm{OK}^+)$ and $(\mathrm{OK}^-)$ of Lemma~\ref{lem:first-direction} fail, and
verify directly that $S$ is not crooked.
\end{exercise}

\begin{exercise}
In Lemma~\ref{lem:up-phase}, explain why the minimum $t=\min(v_+,v_-)$ cannot be replaced by $\max(v_+,v_-)$.
\end{exercise}

\begin{exercise}
Let $\ell,h$ be adjacent indices with $\lo(h)=\lo(\ell)+1$. Show that $\hi(h)=\lo(h)$ and $\hi(\ell)\in\set{\lo(\ell),\lo(\ell)+1}$. Give an
example with chains of open intervals in $\R$ in which $\hi(\ell)=\lo(\ell)+1$.
\end{exercise}

\begin{exercise}
Where exactly is the hypothesis $k\ge6$ used in the proof of Proposition~\ref{prop:splice}? What goes wrong with the middle schedule
$g$ for $k=5$?
\end{exercise}

\begin{exercise}
Deduce from Theorem~\ref{thm:folding} the single-link version: if $r\le s$ and some $E_h$, $E_{h'}$ with $r\le h,h'\le s$ meet $D_0$ and
$D_{k-1}$, then there is a block map with small blocks sending $r$ to the first and $s$ to the last block. (Hint: use the pairs
$(r,r+1)$ and $(s-1,s)$; treat $s\le r+1$ separately.)
\end{exercise}
